\section*{अध्याय \ref{ch:b2:chemical-potential} --- रासायनिक विभव और मिश्रण}

\begin{solution}{exo:b2:chemical-potential:1}
$x_2 = 0.4$ पर स्पर्शरेखा किनारों को $\bar V_1 \approx 14$ और $\bar
V_2 \approx 49$ (मॉडल इकाइयाँ) पर काटती है। तब $0.6 \times 14 + 0.4 \times 49 = 28.0$,
जो $x_2 = 0.4$ पर वक्र से पढ़ा गया मोलर आयतन है, और शुद्ध आयतनों के
$0.6 \times 18 + 0.4 \times 58 = 34.0$ से कम है।
\end{solution}

\begin{solution}{exo:b2:chemical-potential:2}
$n(\text{जल}) = 1000/18.015 = \qty{55.5}{mol}$, $x_1 = 55.5/56.5 = 0.982$;
$p = 0.982 \times 3.17 = \qty{3.11}{kPa}$: 1.8\,\% का अवनमन।
\end{solution}

\begin{solution}{exo:b2:chemical-potential:3}
$p(\ce{O2}) = 0.2095 \times \qty{1.013e5}{Pa} = \qty{2.12e4}{Pa}$;
$c = 1.3 \times 10^{-5} \times 2.12 \times 10^4 = \qty{0.28}{mol/m^3} =
\qty{2.8e-4}{mol/L}$, अर्थात् $2.8 \times 10^{-4} \times 32.0 \times 10^3 =
\qty{8.8}{mg/L}$।
\end{solution}

\begin{solution}{exo:b2:chemical-potential:4}
$Q = (1.0)^2/[(1.0)(0.010)^3] = 1.0 \times 10^6$; $RT\ln Q = 2.479 \times 13.82 =
\qty{34.2}{kJ/mol}$; $\Delta_r G = -32.8 + 34.2 = \qty{+1.4}{kJ/mol} > 0$: हाइड्रोजन में
इतने निर्धन इस मिश्रण में अमोनिया विघटित होती है।
\end{solution}

\begin{solution}{exo:b2:chemical-potential:5}
गिब्स--ड्यूहेम: $x_1\dd\bar V_1 + x_2\dd\bar V_2 = 0$, अतः $\dd\bar V_2 =
-(0.70/0.30) \times (-0.30) = \qty{+0.70}{mL/mol}$।
\end{solution}

\begin{solution}{exo:b2:chemical-potential:6}
$c(\text{कण}) = 2 \times 9.0/58.44 = \qty{0.308}{mol/L} =
\qty{308}{mol/m^3}$; $\Pi = 308 \times 8.314 \times 310 = \qty{7.9e5}{Pa}$,
लगभग \qty{7.9}{bar}। लाल कोशिका के भीतर का \omterm{def:b2:chemical-potential:osmotic-pressure}{परासरण दाब} भी यही है:
जल का कोई कुल प्रवाह नहीं। शुद्ध जल में जल कोशिका में प्रवेश करता है, जो फूलती है और
फट जाती है।
\end{solution}

\begin{solution}{exo:b2:chemical-potential:7}
$c = \Pi/RT = 825/(8.314 \times 298.15) = \qty{0.333}{mol/m^3}$;
\qty{1.000e-4}{m^3} में $n = \qty{3.33e-5}{mol}$, $M = 2.00/3.33 \times 10^{-5}
= \qty{6.0e4}{g/mol}$ (\qty{60}{kg/mol})। $\Delta T_f = 1.86 \times 3.33
\times 10^{-5}/0.100 = \qty{6.2e-4}{K}$: अमापनीय, जबकि \omterm{def:b2:chemical-potential:osmotic-pressure}{परासरण
दाब} जल का \qty{8}{cm} स्तंभ है।
\end{solution}

\begin{solution}{exo:b2:chemical-potential:8}
(a) $I = \qty{1.0e-3}{mol/kg}$, $\log\gamma_\pm = -0.51 \times 0.0316$,
$\gamma_\pm = 0.964$। (b) $I = \frac12(4 \times 1.0 + 1 \times 2.0) \times
10^{-3} = \qty{3.0e-3}{mol/kg}$, $\log\gamma_\pm = -0.51 \times 2 \times 0.0548$,
$\gamma_\pm = 0.879$। (c) $I = \frac12(4 + 4) \times 10^{-3} = \qty{4.0e-3}{mol/kg}$,
$\log\gamma_\pm = -0.51 \times 4 \times 0.0632$, $\gamma_\pm = 0.74$।
\end{solution}

\begin{solution}{exo:b2:chemical-potential:9}
$b = 0.310/1.86 = \qty{0.167}{mol/kg}$; $n = 0.167 \times 0.0500 =
\qty{8.33e-3}{mol}$; $M = 1.50/8.33 \times 10^{-3} = \qty{180}{g/mol}$ (उदाहरण के लिए,
कोई हेक्सोस शर्करा)।
\end{solution}

\begin{solution}{exo:b2:chemical-potential:10}
$x_1\,\dd(Ax_2^2) + x_2\,\dd(Ax_1^2) = 2Ax_1x_2(\dd x_2 + \dd x_1) = 0$ क्योंकि
$x_1 + x_2 = 1$। जब $x_1 \to 0$, तो $\gamma_1 \to \mathrm e^{A}$ और $p_1 =
\gamma_1x_1p_1^* \approx \mathrm e^{A}p_1^*\,x_1$: $k_{H,1} = p_1^*\mathrm e^A$।
$A = 1.1$ के लिए $\mathrm e^{1.1} = 3.0$: हेनरी की प्रवणता राउल्ट की प्रवणता की तिगुनी है।
\end{solution}

\begin{solution}{exo:b2:chemical-potential:11}
(a) वाष्पशील विलेय अपना स्वयं का वाष्प दाब जोड़ता है; वाष्प अब
शुद्ध विलायक नहीं रहती और व्युत्पत्ति (विलयन के साथ साम्य में विलायक की शुद्ध
वाष्प) विफल हो जाती है। (b) यदि विलेय ठोस में प्रवेश करता है, तो
ठोस का \omterm{def:b2:chemical-potential:chemical-potential}{रासायनिक विभव} भी $RT\ln x_1^{\text{s}}$ से घटता है। यदि
ठोस विलयन का संघटन द्रव जैसा ही हो, तो दोनों रेखाएँ समान मात्रा में नीचे
आती हैं और हिमांक नहीं बदलता: अवनमन के लिए ऐसा विलेय
चाहिए जिसे क्रिस्टल अस्वीकार करे।
\end{solution}

\begin{solution}{exo:b2:chemical-potential:12}
हिमांकमिति: $b_{\min} = 0.01/1.86 = \qty{5.4e-3}{mol/kg}$; क्वथनांकमिति:
$0.01/0.513 = \qty{1.9e-2}{mol/kg}$। $\qty{5.4e-3}{mol/L} =
\qty{5.4}{mol/m^3}$ पर $\Pi = 5.4 \times 8.314 \times 298 = \qty{1.3e4}{Pa}$,
जल के एक मीटर से अधिक; परासरणमापी कुछ पास्कल पढ़ लेता है। परासरणमिति
कहीं अधिक सुग्राही है, इसीलिए यह बहुलकों के बड़े मोलर द्रव्यमान मापती है,
जिनकी मोललताएँ अत्यल्प होती हैं।
\end{solution}

\begin{solution}{pb:b2:chemical-potential:1}
\textbf{1.} $\mu_1 = \mu_1^*(T, p) + RT\ln x_1$।
\textbf{2.} डाइऑल $30/62.07 = \qty{0.483}{mol}$, जल $70/18.015 =
\qty{3.886}{mol}$; $x_1 = 3.886/4.369 = 0.889$।
\textbf{3.} $p = 0.889 \times 3.17 = \qty{2.82}{kPa}$।
\textbf{4.} $b = 0.483/0.070 = \qty{6.90}{mol/kg}$।
\textbf{5.} डाइऑल \qty{197}{\celsius} पर उबलता है: कमरे के ताप पर उसका वाष्प
दाब जल के वाष्प दाब से बहुत कम है।
\textbf{6.} $\mu_1^{\text{s}}(T_f) = \mu_1^{*\text{l}}(T_f) + RT_f\ln x_1$।
\textbf{7.} $\ln x_1 = -\Delta_{\text{गलन}}G(T)/(RT)$, जहाँ
$\Delta_{\text{गलन}}G = \mu_1^{*\text{l}} - \mu_1^{\text{s}}$, जो $T_f^*$ पर शून्य है;
गिब्स--हेल्महोल्ट्ज़, $\dd(\Delta_{\text{गलन}}G/T)/\dd T = -\Delta_{\text{गलन}}H/T^2$,
का $T_f^*$ से $T_f$ तक समाकलन $\Delta_{\text{गलन}}G(T_f)/T_f =
\Delta_{\text{गलन}}H(1/T_f - 1/T_f^*)$ देता है, जिससे संबंध मिलता है।
\textbf{8.} $\ln x_1 \approx -x_2 \approx -bM_1$ और $1/T_f - 1/T_f^* \approx
-\Delta T_f/T_f^{*2}$: $\Delta T_f = (RT_f^{*2}M_1/\Delta_{\text{गलन}}H)\,b$।
\textbf{9.} $333.4 \times 18.015 = \qty{6007}{J/mol}$; $K_f = 8.314 \times
273.15^2 \times 0.018015/6007 = \qty{1.86}{K.kg/mol}$।
\textbf{10.} यह कि $\Delta_{\text{गलन}}H$, $T_f$ और $T_f^*$ के बीच नहीं
बदलता (वह बदलता है, किरखॉफ़ से, यहाँ प्रति \qty{3}{K} लगभग 2\,\%)।
\textbf{11.} $\Delta T_f = 1.86 \times 6.90 = \qty{12.8}{K}$: \qty{-12.8}{\celsius}।
\textbf{12.} $1/T_f = 1/273.15 - (8.314/6007)\ln 0.889$: $T_f = \qty{261.6}{K}$,
\qty{-11.6}{\celsius}।
\textbf{13.} यथातथ आदर्श संबंध पर, जो तनु विलयन नहीं मानता (यहाँ 11\,\% अणु
डाइऑल हैं)। दोनों \omterm{def:b2:chemical-potential:ideal-mixture}{आदर्श मिश्रण} और स्थिर गलन एन्थैल्पी मानते हैं;
हाइड्रोजन आबंधों से जुड़े जल और डाइऑल आदर्श नहीं हैं, इसलिए वास्तविक हिमांक
भिन्न होता है।
\textbf{14.} शुद्ध बर्फ़: डाइऑल द्रव में ही रहता है।
\textbf{15.} बर्फ़ के रूप में जल हटने से द्रव में $x_1$ घटता है, और
संबंध के अनुसार साम्य ताप गिरता है: शीतलक एक बिंदु पर नहीं, बल्कि तापों के एक
परिसर में जमता है।
\textbf{16.} $K_b = 8.314 \times 373.12^2 \times 0.018015/40\,650 =
\qty{0.513}{K.kg/mol}$।
\textbf{17.} $\Delta T_b = 0.513 \times 6.90 = \qty{3.5}{K}$।
\textbf{18.} $1/T = 1/373.12 - (8.314/40\,650)\ln(2.0/1.013)$: $T = \qty{393.5}{K}$,
लगभग \qty{120}{\celsius}।
\textbf{19.} ढक्कन: \qty{20}{K}, जबकि डाइऑल के लिए \qty{3.5}{K}।
\textbf{20.} $x_1 = \exp[-(6007/8.314)(1/253.15 - 1/273.15)] = 0.811$।
\textbf{21.} $x_2 = 0.189$: $w = 0.189 \times 62.07/(0.189 \times 62.07 + 0.811
\times 18.015) = 0.44$।
\textbf{22.} तनु नियम: $b = 20/1.86 = \qty{10.75}{mol/kg}$, $w = 10.75 \times
62.07/(1000 + 10.75 \times 62.07) = 0.40$। रैखिकीकरण ($\ln x_1 \approx
-x_2$, $x_2 \approx bM_1$) इतनी उच्च सांद्रता पर डाइऑल के प्रति मोल अवनमन को
बढ़ा-चढ़ाकर आँकते हैं, इसलिए तनु नियम कम डाइऑल माँगता है।
\textbf{23.} आदर्श-मिश्रण मॉडल में डाइऑल के द्रव्यमान अंश
$\boldsymbol{w \approx 0.44}$ वाला शीतलक \qty{-20}{\celsius} तक द्रव रहता है।
\end{solution}
